\section{Implementation inventory: learning, representation, and readout}
\label{sec:implementation_inventory}
This inventory covers retained historical protocols; the current direct-field
main experiments are documented in Section~\ref{sec:direct_main}.
The simulation evidence has three distinct parts: an estimate changes in
response to observed incompatibility; that estimate represents evaluated
partials; and a readout from the estimate supplies a task command.
Tables~\ref{tab:implementation_main} and
\ref{tab:implementation_support} identify the implemented mechanisms rather
than treating every example as a run of one optimizer. The periodic-plant
experiment in Section~\ref{sec:normalized_fourier_results} integrates the
coupled normalized law. The retained contact and Burgers protocols implement the separately
specified constrained coefficient extension. Their results support learning,
representation, and readout within the stated families, without establishing
that their updates satisfy the normalized-flow convergence hypotheses.
As distinguished in Section~\ref{sec:scope}, these are retained
coefficient-family demonstrations, not the required all-entry direct-field
experiments.

\subsection{Coefficient-induced partials versus directly learned entries}
For a prescribed feature family with all supplied terms retained in the
implicit relation, write
\begin{equation}
 f_\theta(z)=Q_\theta\zeta(z),\qquad
 F_\theta(z)=Q_\theta D_z\zeta(z).
 \label{eq:inventory_representation}
\end{equation}
At a fixed context, a parameter update induces the exact change
\begin{equation}
 \Delta_{\rm learn}F(z)
 =(Q_{\theta^+}-Q_\theta)D_z\zeta(z).
 \label{eq:inventory_partial_update}
\end{equation}
Supplied coefficient entries have zero learning increment. Thus the code
learns residual coefficients and evaluates their induced partials; it does
not directly adapt every partial entry with unknown algebraic structure.
The chain-rule identity does not make these two learning mechanisms
equivalent. Context motion
can also change the evaluated partials even when the coefficients are frozen.
Neither a rapidly varying curve nor coefficient convergence alone establishes
accurate physical partials. The error diagnostics compare estimates and
evaluator truth at the same context, using a declared scale convention.

For the seven-coefficient normalized periodic family, this relation gives
the explicit continuous learning contribution
\begin{equation}
 \left.\partial_t F_q(z)\right|_z
 =-\eta\bigl[qR-(qRq^{\mathsf T})q\bigr]S D_z\zeta(z).
 \label{eq:inventory_normalized_partials}
\end{equation}
Its $F_q$ is (\ref{eq:nf_jacobian}), including both clock-product terms.
Every candidate is an exact Jacobian, and all seven coefficients, including
the input coefficient, adapt. Dividing the displayed partials by that
learned input coefficient is a reporting convention, not a fixed training
coefficient. The generic six-coefficient direct fields in
Sections~\ref{sec:unanchored} and~\ref{sec:direct_lift_results} instead have
$\Delta H(z)=\Delta q\,\Phi(z)$ for their field basis $\Phi$; they need not
be Jacobians and are not covered by (\ref{eq:inventory_representation}).

\begin{table*}[t]
\caption{Retained historical simulation mechanisms. Paths in this table are relative to
\texttt{experiments/}, except the explicitly identified legacy Burgers solver.
The historical normalized Fourier study uses its own controlled trajectory;
replay results are identified separately in
Section~\ref{sec:normalized_fourier_results}.}
\label{tab:implementation_main}
\centering
\setlength{\tabcolsep}{4pt}
\renewcommand{\arraystretch}{1.12}
\begin{tabular}{@{}p{.14\textwidth}p{.29\textwidth}p{.34\textwidth}>{\raggedright\arraybackslash}p{.17\textwidth}@{}}
\toprule
Study & Observations and supplied information & Learned representation and actual update & Simulation and figure code\\
\midrule
Normalized Fourier & Complete-feature secants of $(hr_m,hx,-u)$.
State, input, phase $\phi=4t$, harmonic family, scales, and zero residual
level are supplied; rate measurement is noisy. &
Seven unit-row coefficients $q$; $a=qS$.
Coupled $\dot R=ww^{\mathsf T}-R/8$ and
$\dot q=-\eta[qR-(qRq^{\mathsf T})q]$ are numerically integrated with
changing $R$. The full $F_q$ in (\ref{eq:nf_jacobian}) determines the
learned relation used for command readout and current-rate completion.
No partial column is fixed numerically, but their algebraic dependence
remains prescribed by the residual family.
& \path{normalized_fourier.py};
\path{plot_normalized_fourier_main.py};
\path{report_normalized_fourier.py}.\\[3pt]
Contact & Material integral windows from position and noisy force;
separate measured velocity residuals $v_m-Bu$.
The basis, $\rho=4$, unit rate/velocity coefficients, coefficient boxes,
and zero levels are supplied. Force direction informs the control readout.
& Six material coefficients $(s,a,r,b)$ and four entries of $B$.
Each fit uses 48 projected-gradient steps
(\ref{eq:audit_exact_update}). Material partials are
(\ref{eq:audit_contact_partials}); the kinematic input block is $-B$.
These are integral/level fits, not raw-context tangent normalized flow.
The learned material and actuation estimates enter the contact readout.
& \path{contact.py};
\path{common.py};
\path{make_figures.py}.
Transfer: \path{contact_excitation.py},
\path{plot_excitation.py}.\\[3pt]
Burgers & Spatially averaged complete-feature secants of
$(w_{t,m},ww_s,-w_{ss},-u)$.
Exact state and spectral spatial channels, full actuation, unit $w_t$
coefficient, bounds, and zero level are supplied. &
Three coefficients $(c_a,\nu,b)$; eight scaled projected-gradient steps
(\ref{eq:audit_exact_update}).
$F=(c_aw_s,c_aw,-\nu,-b,1)$.
Both partials of the nonlinear product are represented.
The three learned coefficients enter the inverse control readout;
the discretized family and spatial operators remain prescribed.
& \path{burgers_online.py};
\path{make_figures.py}.
Solver: \path{examples/legacy_suite/simulations/burgers.py}
(repository relative).\\
\bottomrule
\end{tabular}
\end{table*}

\subsection{Updates used by the supporting comparisons}
For the constrained coefficient studies, let $\theta$ denote scaled free
coefficients and $\mathcal W$ their box. Their current-candidate objective
is $J_n(\theta)=\tfrac12\theta^{\mathsf T}H_n\theta+g_n^{\mathsf T}\theta+
\mathrm{const}$, with supplied terms retained in the full implicit row.
The exact memory and projected update are given in
(\ref{eq:audit_exact_update}). The matched EWLS comparison minimizes this
same box-constrained objective; it adds no ridge penalty. The instantaneous
comparison instead takes
\begin{equation}
 \theta^+=\Pi_{\mathcal W}\!\left[
 \theta-\frac{a_n(a_n^{\mathsf T}\theta+b_n)}
 {\max(\|a_n\|^2,10^{-12})}\right],
 \label{eq:inventory_nlms}
\end{equation}
where $(a_n,b_n)$ is the current scaled observed implicit feature channel.
Retaining a covariance for diagnostics does not make this NLMS update a
memory-based estimator.

The Acrobot step is
\begin{equation}
 \theta^+=(I+\eta hH_n)^{-1}(\theta-\eta h g_n).
 \label{eq:inventory_proximal}
\end{equation}
Its memory uses the actual sample weight $h$:
$H_n=\beta H_{n-1}+hV_n^{\mathsf T}V_n$ and
$g_n=\beta g_{n-1}+hV_n^{\mathsf T}k_n$, where $V_n\theta+k_n$
is the complete-feature secant. It is a proximal step, without a persistent
ridge penalty or a positivity projection.

In the active-flow prototype the objective instead contains a ridge term:
\begin{align}
 \theta_n^\star&\in\arg\min_{\theta\in\mathcal W}
 \left[J_n(\theta)+\frac{\lambda_n}{2}
             \|\theta-\theta_0\|^2\right],\\
 \theta^+&=.75\theta+.25\theta_n^\star,
 \label{eq:inventory_flow_relaxation}
\end{align}
with $\lambda_n=10^{-7}\max(\tr G_n/3,10^{-10})$ and $G_n$ the
unregularized feature Gram matrix. These distinctions affect the optimized
criterion and cannot be omitted when relating the experiments to theory.

\begin{table*}[t]
\caption{Supporting mechanisms and source files. In the code column,
E denotes \texttt{experiments/}; L denotes
\texttt{examples/legacy\_suite/simulations/}. The cited sections specify
configurations, records, and metric conventions. A plotting or diagnostic
script does not define an additional learning rule.}
\label{tab:implementation_support}
\centering
\setlength{\tabcolsep}{4pt}
\renewcommand{\arraystretch}{1.10}
\begin{tabular}{@{}p{.14\textwidth}p{.29\textwidth}p{.34\textwidth}>{\raggedright\arraybackslash}p{.17\textwidth}@{}}
\toprule
Study & Observations and adaptive/supplied quantities & Actual update, representation, and readout status & Source / postprocessor\\
\midrule
Original Fourier and estimator comparisons & Same complete features as
the periodic plant; six physical harmonic coefficients adapt, while the
input coefficient is supplied. Phase, family, boxes, and zero level are known.
The $K=0$ ablation learns only two constants. &
Eight projected-gradient steps, box EWLS minimization, or
(\ref{eq:inventory_nlms}), as identified in
Sections~\ref{sec:fourier_protocol} and~\ref{sec:estimator_comparison}.
$F$ is (\ref{eq:audit_fourier_partials}). Same-stream shadows are
distinguished from methods controlling their own trajectories. &
E: \path{fourier_online.py},
\path{estimator_comparison.py};
\path{make_figures.py},
\path{plot_estimator_comparison.py}.\\[3pt]
Acrobot & Complete-feature secants from known positions, velocities and
commands, with noisy acceleration. Five mechanical coefficients adapt;
the actuation column $B=(0,1)^{\mathsf T}$ and mechanical family are supplied.
& Proximal step (\ref{eq:inventory_proximal}).
$F=D_z[M_c(q)\alpha+h_c(q,\dot q)-Bu]$ at fixed $c$.
Identification precedes a separate local LQR episode with the final
coefficients frozen; invalid final models are rejected for readout.
See Figure~\ref{fig:acrobot_archived}. &
L: \path{acrobot.py} generates both records and figures.\\[3pt]
Active-flow prototype & Noiseless full velocity fields and held command;
midpoint fields and secant acceleration form complete-feature differences.
Three coefficients adapt. Spatial operators, geometry, forcing shape,
and acceleration identity are supplied. &
Ridge box solve with 25\% relaxation,
(\ref{eq:inventory_flow_relaxation}).
$F=[c_aDN-\nu L-DK,\ I,\ -b_gB]$.
The adaptive branch continues learning during control; the frozen branch
does not. Differential error probes a specified direction, not a full
operator norm (Section~\ref{sec:afc_prototype}). &
L: \path{airfoil.py} generates records, parameter and control plots.
E: \path{plot_active_flow_fields.py} renders the saved velocity as
vorticity and instantaneous streamlines.\\[3pt]
Reconstructed unanchored Fourier & 4096 noiseless phase/rate samples form
one fixed batch covariance $R$. All six unit-row coefficients adapt;
the harmonic field family and phase are supplied. &
Retracted step $q^+=\operatorname{unit}(q-\alpha g)$,
$g=qR-(qRq^{\mathsf T})q$, with energy/phase-guard backtracking.
$H_q$ need not be a Jacobian. This is a field-learning diagnostic,
without a simulated control episode
(Section~\ref{sec:unanchored}). &
L: \path{fourier.py} generates the checked reconstruction and its figure.\\[3pt]
Direct field and exact-Jacobian companion & Exact phase and one noisy
rate channel; all six field coefficients or all four Jacobian-family
coefficients adapt. The basis is supplied; no derivative entry is fixed.
& $R^+=\beta R+(1-\beta)\psi\psi^{\mathsf T}$, then
$q^+=\operatorname{unit}(qe^{-\eta hR^+})$.
This is a frozen-memory normalized substep.
The six-dimensional $H_q$ need not be integrable; the companion has
$F_q=(q_0+q_2\cos x-q_3\sin x,q_1)$.
Neither study has a controlled plant (Section~\ref{sec:direct_lift_results}).
& E: \path{direct_lift_audit.py},
\path{direct_jacobian_audit.py};
\path{plot_direct_lift_audit.py},
\path{audit_direct_jacobian.py}.\\[3pt]
Additional diagnostics & Existing learned and comparison records;
no new observations or adaptive states. The contact transfer study
uses the contact learner before freezing it for transfer. &
Figure~\ref{fig:unknown_partial_tracking} recalculates error on learned
entries only. Extended plots retain their underlying learner.
Freezing coefficients during contact transfer does not freeze the
context-dependent partials (Section~\ref{sec:excitation}). &
E: \path{derivative_tracking_audit.py},
\path{make_figures.py},
\path{plot_excitation.py}.
\path{repair_supplement_layout.py} changes layout only.\\
\bottomrule
\end{tabular}
\end{table*}

Here $\operatorname{unit}(v)=v/\|v\|$ for a nonzero row.
The frozen-memory exponential step in Table~\ref{tab:implementation_support}
is distinct from numerical integration with changing $R$ in
(\ref{eq:nf_coupled}). Its six-dimensional evaluated field is
\begin{equation}
 H_q(x)=\bigl(q_0+q_2\cos x+q_4\sin x,
                q_1+q_3\cos x+q_5\sin x\bigr).
\end{equation}
Generic $q_3,q_5$ violate the cross-partial condition for a scalar
potential in $(x,y)$. The four-coefficient companion is consequently a
separate exact-Jacobian test, not a relabeling of that field experiment.
The normalized periodic study also includes a bottom-eigenvector minimizer
of the same unit-row memory objective and a frozen prior. They are
same-stream estimators, not independently executed controllers.

The circle and mechanical constructions in
Section~\ref{sec:constructions} are analytic examples without simulated
learning. The ellipse discussion is likewise retained only as a mathematical
construction: its original implementation and raw records are unavailable,
so it supplies no empirical learning or completion result. The unanchored
Fourier figure now uses the checked executable reconstruction; it does not
claim recovery of the earlier historical figure's unknown lineage.

\subsection{What the simulation chain demonstrates}
The retained historical simulations evaluate pre-update incompatibility, estimated versus
true evaluated partials, and task performance separately. The normalized
periodic learner controls its own simulated plant while all seven
coefficients adapt. The same saved pre-update representation also supports
the different query
\begin{equation}
 \widehat r=\frac{a_u u-(a_\kappa h)x}{a_\rho h},
 \label{eq:inventory_rate_readout}
\end{equation}
using the observed state, input, and phase, without taking the measured
rate as a readout input or retraining the representation. However,
substituting the command (\ref{eq:nf_control}) gives
$\widehat r=r_d$ whenever both readouts are admissible and the command
guard is inactive. Its error therefore measures how closely the actual
plant rate achieves that requested rate. It is a readout-consistency
diagnostic on the same controlled trajectory, not an independent
prediction experiment. The two formulas show algebraic reuse of one
learned relation for two choices of computed channel; they do not supply
two independent empirical validations. The completion denominator is
checked separately. No future prediction or held-out trajectory result
is claimed in these historical records. The retained contact and Burgers
protocols similarly use their current learned
structured relations in active control, with their distinct observation
mechanisms and supplied physical terms disclosed. This is numerical evidence
for coefficient learning and readout within the specified implicit families.
It does not establish direct learning of all partial entries without the
known physical algebra, even when the resulting partial estimates are
accurate. That distinction is separate from whether the numerical optimizer
implements its stated law correctly.

True coefficients and clean evaluator partials are not derivative targets
for the adaptive branches. The plant uses truth to generate physical
responses, and explicitly named oracle comparisons use it by design.
The causal sequence and information checks for these retained protocols are
documented in the respective sections. Calibrated zero levels remain
supplied for the control readouts; learning a differential alone cannot
recover an additive level. A correct feature family, adequate information,
and an admissible readout block remain substantive assumptions.

Finally, a shadow estimator's partial accuracy is evidence about its
representation on the generating controller's data, not about a controller
that was never run. Offline identification followed by frozen-model control
is also distinguished from simultaneous adaptation and control. These
separations delimit what each retained simulation demonstrates. The
separately identified direct-field learning, representation, and readout
experiments in Section~\ref{sec:direct_main} supply the current main evidence;
they do not reinterpret these historical coefficient fits as direct learning.
